\begingroup
\emergencystretch=3em
\section{Выводы}
В работе рассмотрена многопоточная обработка матрицы вещественных
чисел фильтрами эрозии и наращивания. Реализовано независимое
\(K\)-кратное применение каждого фильтра с обрезанием окна на границе.
Общее число участников ограничивается параметром запуска;
главный поток также выполняет часть вычислений.

Работа с потоками вынесена в интерфейс \texttt{mt}, а Linux-операции
создания, ожидания и синхронизации находятся в платформенной реализации.
Разбиение по непересекающимся диапазонам строк позволяет выполнять
вычисления без мьютекса на каждом элементе. Двойная буферизация
и барьер между проходами обеспечивают чтение полностью рассчитанного
предыдущего состояния матрицы.

Стартовая блокировка предотвращает зависание на барьере при частичном
создании группы потоков. После ошибки уже созданные участники выходят
до начала вычислений, а главный поток ожидает их завершения.

В ранее подготовленной контрольной Linux-сборке проверены эталонная
матрица, граничные параметры и количество потоков средствами ОС.
Измерения производительности текущей версии выполнены отдельно
на компьютере студента в WSL2 с процессором Intel Core i7-13650HX.

Для матрицы \(256\times256\) наименьшая медиана времени получена
при двух потоках: ускорение 1.692 и эффективность 84.6\%.
Для матрицы \(1024\times1024\) наименьшая медиана получена
при восьми потоках: ускорение 3.916 и эффективность 49.0\%.
Увеличение объёма работы улучшило использование нескольких потоков.
Рост числа потоков не обеспечил пропорционального ускорения.

В расчётах использованы медианы пяти повторов каждой конфигурации.
Разброс отдельных запусков сохранён в исходных CSV; результаты
характеризуют конкретную сборку и условия измерения, а не гарантированную
производительность на другом компьютере.
\clearpage
\endgroup
